Physics-informed neural networks (PINNs) provide cheap, mesh-free,
differentiable surrogates for partial differential equations, but their
pointwise accuracy is known to be \emph{spatially non-uniform}: a small
subdomain typically carries a disproportionate share of the error. Existing
remedies---adaptive sampling, loss reweighting, causal training, domain
decomposition---rebalance the optimisation but still return a single global
predictor that the user must trust everywhere. We argue instead for
\emph{learning with abstention}: a frozen PINN is paired with a classical
solver, and a learned rejector decides at each spatial location whether to
commit to the PINN prediction or to abstain to the solver. We cast this as
a regression abstention problem with an input-dependent cost derived from
the PDE residual and a linearised error-transport estimate, derive its
Bayes-optimal rejector, and design a consistent surrogate for the
predictor--rejector formulation extended from a categorical predictor with
$0$--$1$ loss to a continuous predictor with a non-negative discrepancy.
We prove an $\Hcal$-consistency bound for the proposed surrogate and show
that the induced hybrid solution inherits a controllable PDE error from
the rejection decisions.
